% The sensor's buffer: one tag byte and six data bytes per entry.
\begin{tikzpicture}[font=\sffamily\scriptsize]

% one entry: tag cell plus data cell
\newcommand{\fifoentry}[4]{%
  \node[draw=linecol, fill=#3, anchor=south west, minimum width=9mm,
        minimum height=6mm, inner sep=1pt, font=\sffamily\tiny] at (0mm,#1mm) {#2};
  \node[draw=linecol, fill=white, anchor=south west, minimum width=32mm,
        text width=32mm, align=center, minimum height=6mm, inner sep=1pt,
        font=\sffamily\tiny] at (9mm,#1mm) {#4};}

\node[font=\sffamily\bfseries\small, anchor=south west] at (0mm,62mm) {the sensor buffer};
\node[font=\sffamily\tiny, anchor=south west] at (0mm,57mm) {tag};
\node[font=\sffamily\tiny, anchor=south west] at (9mm,57mm) {six data bytes};

\fifoentry{0}{gyro}{hwcol}{oldest entry, replaced first}
\fifoentry{6}{accel}{kercol}{}
\fifoentry{12}{gyro}{hwcol}{}
\fifoentry{18}{accel}{kercol}{}
\draw[mmgap] (-2mm,26mm) -- (43mm,26mm);
\node[font=\sffamily\tiny, anchor=north west, align=left, text width=44mm] at (58mm,62mm)
  {The break is a break in scale. The depth in entries is in the sensor
   datasheet and is never assumed from a sibling part.};
\fifoentry{30}{accel}{kercol}{}
\fifoentry{36}{gyro}{hwcol}{}
\fifoentry{42}{accel}{kercol}{newest entry}
\draw[deadline] (-2mm,48mm) -- (43mm,48mm);
\node[font=\sffamily\tiny, anchor=south west] at (-2mm,49mm) {watermark level: the interrupt fires here};

% ---- the MCU side
\node[blk, fill=usrcol, text width=34mm, anchor=south west, minimum height=12mm]
  (land) at (58mm,34mm) {landing buffer\\{\tiny 2 kB, the level is clamped to it}};
\node[blk, fill=kercol, text width=34mm, anchor=south west, minimum height=10mm]
  (r1) at (58mm,18mm) {accelerometer ring};
\node[blk, fill=kercol, text width=34mm, anchor=south west, minimum height=10mm]
  (r2) at (58mm,4mm) {gyroscope ring};
\node[blk, fill=fwcol, text width=26mm, anchor=south west, minimum height=10mm]
  (dc) at (100mm,11mm) {drop counters\\{\tiny overrun and tag gap}};

\draw[flow] (43mm,44mm) -- node[lbl, above]{one burst} (land.west |- 0mm,44mm);
\draw[flow] (land.south) ++(-8mm,0) -- ++(0,-4mm) -| (r1.north);
\draw[flow] (land.south) ++(8mm,0) -- ++(0,-18mm) -| (r2.north);
\draw[flow] (r1.east) -- ++(4mm,0) |- (dc.west);

\node[umlnote, text width=52mm, anchor=north west] at (0mm,-4mm)
  {The tag is what lets one buffer hold several streams, and it is what makes
   the drain a decode rather than a copy. It is also the second drop detector:
   a gap in the tag sequence is a lost sample even when the overrun flag was
   missed.};
\node[umlnote, text width=46mm, anchor=north west] at (58mm,-4mm)
  {Continuous mode replaces the oldest entry and raises the overrun flag. The
   stop-when-full mode keeps the oldest and stops advancing, which is much
   harder to notice.};
\end{tikzpicture}
